\chapter{New results on $1/2$-flow-pairs} % chapter* je necislovana kapitola
\label{ch:flow_pairs}

In this chapter we show that $1/2$-flow-pairs are worth further research, since they are in a close relation with several flow concepts. In the first part, we show their existence on snarks of oddness $2$ (and hence also on weak snarks), which form the majority of currently known snarks. We also show, that $1/2$-flow-pair implies a rich $7$-flow, what is an improvement of the upper bound for all snarks of oddness $2$.

\section{Oddness $2$}

One of the first results on snarks admitting a NZ $5$-flow was the class of snarks with oddness two (see Theorem \ref{th:oddness_1D}). The main idea of the proof is adding a particular edge $uv$ such that the resulting graph has a NZ $4$-flow, and then finding an oriented $uv$-path, where the flow values are all positive, so we can get rid of the flow on the edge $uv$.

We generalise this approach and prove a similar theorem for flow-pairs. However, in this setting, we need to find two dart-disjoint $uv$-paths before deleting the edge $uv$.

\begin{theorem}
Each snark $G$ of oddness $2$ admits a $1/2$-flow-pair.
\end{theorem}

\begin{proof}
Consider a $2$-factor $F$ of $G$ with two odd cycles. Then construct an NZ $(\mathbb Z_2\times \mathbb Z_2)$-flow $\varphi$ on $G$ as follows – assign $(0,1)$ to edges in $E(G)\setminus F$ and, on the $2$-factor $F$, alternate flow values $(1,1)$ and $(1,0)$. Since all elements of $\mathbb Z_2\times \mathbb Z_2$ are involutions, assigning a value to an edge is equivalent to assigning the same value to both corresponding darts. On both odd cycles, there is one conflicting vertex, where the sum of flow values is $(0,1)$. Let these conflicting vertices be $u$ and $v$. Join these two vertices by a new edge $e$ with an assigned flow value $(1,1)$ and denote the new graph by $G'$. Thus, we have an NZ $(\mathbb Z_2\times \mathbb Z_2)$-flow $\varphi$ on $G'$. By Theorem \ref{th:finite_group_flows} we transform this NZ $(\mathbb Z_2\times\mathbb Z_2)$-flow to a $0/1$-flow-pair. However, we know that $\varphi\left(\vv{uv}\right)$ is either $(0, 1)$ or $(0, -1)$. In the latter case, we can multiply the second coordinate of the whole flow by $-1$. Therefore we can WLOG assume $\varphi\left(\vv{uv}\right)=(0,1)$.

Multiply the second coordinate of $\varphi$ by two to obtain $\varphi'=(\varphi'_1,\varphi'_2)=(\varphi_1, 2\varphi_2)$. Observe that the flow values of $\varphi'$ are feasible for a $1/2$-flow-pair and $\varphi'\left(\vv{uv}\right)=(0,2)$. Call a dart $d$ \emph{good} if the value of $\varphi'\left(d\right) + (0,1)$ is still a feasible flow value of the $1/2$-flow-pair. Now, if we find two oriented cycles containing the dart $\vv{vu}$ such that all other darts in cycles are good, we can get rid of the flow on the edge $uv$ and then delete it. Clearly, we cannot take the same dart twice, but when we take two darts corresponding to the same edge in $G$ and send $(0,1)$ through each of them, the resulting flow value on $e$ is the same as before the process, so this is no issue. Now, we construct a graph $H$ of good darts, formally $V(H)=V(G)$ and
$$E(H)=\{\vv{xy}\mid xy\in E(G) \land (\varphi_1'(\vv{xy})=\pm 1 \lor \varphi_2'(\vv{xy})=2)\}.$$
Now we aim to find two dart-disjoint $uv$-paths in $H$. From Theorem \ref{th:oriented_menger}, there exist two dart-disjoint $\vv{uv}$-paths in $H$ if and only if each $uv$-dart-separator of $H$ contains at least $2$ darts from the component containing $u$ to the component containing $v$. So it remains to prove this condition.

Consider any $uv$-dart-separator $\gamma_{G'}$ of $G'$. From the generalised flow conservation, $\sum_{d\in\gamma_{G'}} \varphi'\left(d\right)=(0,0)$. If there is a dart in $\gamma_{G'}$ having a flow value with non-zero first coordinate, there must be at least one more such dart due to the sum of flow values. For both of them, there exists a corresponding dart in $H$. Hence, we have at least two darts in the corresponding $uv$-dart-separator $\gamma_H$ of $H$. In the remaining case, only possible flow values on the darts of $\gamma_{G'}$ are $(0,\pm 2)$.

Observe that the number of darts with flow value $(0, 2)$ must be the same as the number of darts with flow value $(0, -2)$ due to the zero sum. This also argues that $|\gamma_{G'}|$ is an even number. One of the $(0,2)$ darts is the new $uv$ edge, and for the others, there is a corresponding dart in $\gamma_H$. Hence, if $|\gamma_{G'}|>4$, there are at least three $(0,2)$ darts in $\gamma_{G'}$, and so at least two $(0,2)$ darts in $\gamma_H$. If $|\gamma_{G'}|=2$, there is a bridge in $G$, which contradicts having any NZ flow. So we are left with $|\gamma_{G'}|=4$.

From $|\gamma_{G'}|=4$ we already know that $|\gamma_{G}|=3$. Since $G$ as a snark is cyclically $4$-edge-connected, at least one of the components separated by $\gamma_{G}$ must be an acyclic $3$-pole with $k$ inner vertices. Let $k$ denote the number of its inner vertices of degree $3$. We will use the fact that this $3$-pole is a tree $T$ on $k+3$ vertices with $3$ leaves, and hence $|E(T)|=k+2$. Moreover, from the folklore handshaking lemma, we get
\begin{equation*}
 3k+3 = \sum_{v\in V(T)} \deg(v) = 2|E(T)|=2k+4,
\end{equation*}
and so $k=1$. However, the only cases when the $uv$-dart-separator separates one vertex, are the separators around $u$ and $v$. Both of these separators contain a flow value with non-zero first coordinate, so this case is already solved.
\end{proof}

\begin{corollary}
Each weak snark admits a $1/2$-flow-pair.
\end{corollary}

\section{Rich $7$-flow}

Now we move to the relation of rich flows and flow-pairs. Up to now, the smallest upper bound proved for some infinite family of snarks is in the Proposition \ref{prop:weak_rich_flow}, which shows that if we can delete an edge, smooth the vertices of degree $2$ and obtain a colourable graph, the original graph has the rich flow number at most $9$. However, from the $1/2$-flow-pair, we are able to construct a rich NZ $7$-flow. This means, that for graphs with oddness $2$, we tightened the gap between the construction and the bound from $3$ to only $1$.

\begin{theorem}
Let $G$ be a cubic graph with $1/2$-flow-pair. Then $G$ has a rich nowhere-zero $7$-flow.
\end{theorem}

\begin{proof}
Denote the $1/2$-flow-pair of $G$ as $\varphi=(\varphi_1, \varphi_2)$. First, we construct a circular nowhere-zero $7$-flow $\varphi'$. Then, we will adjust non-integral values of $\varphi'$ to $\varphi''$, which will be the sought rich nowhere-zero $7$-flow.

Let $\varphi'\left(d\right)=9\varphi_1\left(d\right)/2 + \varphi_2\left(d\right)/2$. We see that
\begin{align*}
 \varphi_1\left(d\right)&=-1&&\Rightarrow\varphi'\left(d\right)\in[-6, -3],\\
 \varphi_1\left(d\right)&=0&&\Rightarrow\varphi'\left(d\right)\in\{-3/2, -1, 1, 3/2\},\\
 \varphi_1\left(d\right)&=1&&\Rightarrow\varphi'\left(d\right)\in[3, 6].
\end{align*}
This shows that $\varphi'$ is a circular nowhere-zero $7$-flow.

From Proposition \ref{prop:chebyshev_normal_form}, we can construct an integral nowhere-zero $7$-flow $\varphi''$. Moreover, its flow values satisfy
\begin{align*}
 \varphi_1\left(d\right)&=-1&&\Rightarrow\varphi''\left(d\right)\in\{-6, -5, -4, -3\},\\
 \varphi_1\left(d\right)&=0&&\Rightarrow\varphi''\left(d\right)\in\{-2, -1, 1, 2\},\\
 \varphi_1\left(d\right)&=1&&\Rightarrow\varphi''\left(d\right)\in\{3, 4, 5, 6\}.
\end{align*}
This is the key observation to prove that $\varphi''$ is rich. For the sake of contradiction, assume a vertex $v$, in which two darts are either confluent or contrafluent. They cannot be confluent, since otherwise the third edge would have zero flow value. Hence, denote the common value on the contrafluent darts as $x$ and on the third dart, there is $-2x$, which is in absolute value greater. Hence, the two equal flow values must have the smallest absolute value. From Lemma \ref{lem:flow_pair_2_factor}, each vertex is surrounded by one dart with $\varphi_1\left(d\right)=0$ and two darts with $\varphi_1\left(d\right)=1$. But absolute values of $\varphi''\left(d\right)$ on darts with $\varphi_1\left(d\right)=0$ are strictly smaller than absolute values of $\varphi''\left(d\right)$ on darts with $\varphi_1\left(d\right)\neq 0$. Therefore, there cannot be two darts with the same smallest absolute value, a contradiction. Finally, we get that $\varphi''$ is a rich NZ $7$-flow.
\end{proof}

\begin{corollary}
Each snark $G$ of oddness $2$ has a rich NZ $7$-flow.	
\end{corollary}